% sections/related.tex -- Related work (MMHM v2)
% Citation keys: refs.bib. Verification log: REFERENCES_VERIFICATION.md.
\section{Related work}\label{sec:related}

\paragraph{Static homology and persistence computation.}
Over a field, Betti numbers are ranks of boundary matrices, and the standard
way to obtain them, together with the persistence barcode of a filtration, is
column reduction in a fixed total order of the cells: each column receives
additions from earlier columns until its lowest non-zero entry (its
\emph{low}) differs from that of every other non-zero column
\cite{edelsbrunner2002topological,zomorodian2005computing,edelsbrunner2010computational}.
The worst case is cubic in the number of cells
\cite{edelsbrunner2010computational}, and practical performance comes from a
small set of devices. \emph{Clearing} (the \emph{twist}) processes dimensions
from the top down and sets to zero every column whose cell is already known to
be the low of a reduced column one dimension up \cite{chen2011twist};
\emph{clear and compress} pairs columns locally in chunks and compresses the
remainder before a global reduction \cite{bauer2014clear}. \emph{Apparent
pairs} are decided by a purely local condition on the boundary matrix, need no
column additions, and are both persistence pairs and gradient pairs of a
discrete Morse function; Ripser combines them with clearing, cohomology and an
implicit coboundary \cite{bauer2021ripser}. PHAT \cite{bauer2017phat} and
GUDHI \cite{maria2014gudhi} package these variants, and Otter et
al.~\cite{otter2017roadmap} benchmark the available software. All of this is
static: an edited complex is a new input. Our rebuild step is exactly such a
static reduction with clearing, modified to retain the change of basis~$V$;
our \method{full-*} baselines rerun it, without~$V$, after every update.

\paragraph{Discrete Morse theory and coreductions.}
Forman's discrete Morse theory \cite{forman1998morse,forman2002user} replaces
a complex by a chain-homotopy-equivalent Morse complex on the critical cells
of an acyclic matching. Algebraic Morse theory
\cite{skoldberg2006morse,jollenbeck2009minimal} extends the construction to
chain complexes of free modules and obtains the Morse complex by eliminating the
matched pairs, whose incidences are invertible. Algorithmically, matchings are built by collapses or
coreductions and used to shrink the input before any linear algebra:
coreduction homology \cite{mrozek2009coreduction}, extended to regular CW
complexes \cite{dlotko2011coreduction}, to maps
and persistence \cite{mrozek2010coreduction,harker2014discrete}, and filtered
Morse complexes that preserve persistent homology \cite{mischaikow2013morse}.
On 2-manifolds, matchings with the minimum number of critical cells can be
constructed in linear time \cite{lewiner2003optimal}; in general, Benedetti
and Lutz \cite{benedetti2014random} collapse a free face chosen at random and
declare a cell critical only when no free face is left, and use the resulting
distribution of critical-cell counts to measure how complicated a
triangulation is. Our matching is a deterministic variant of their collapse
strategy. We use it differently from the works
above, because we never form the Morse complex. Its boundary entries are sums
over gradient paths, so an edit to a \emph{matched} cell can change the
boundaries of critical cells far from the edit; this is one of the two defects
that invalidated the localised update of version~1 \cite{ouyang2025mmhm_v1}.
Instead we use the matching only to choose the total order of the
\emph{unreduced} boundary matrix, so that every gradient pair is an apparent
pair. This step is elementary and close to known results: the collapse removes
each paired coface immediately before its face, so the two are adjacent in the
reverse removal order, and a facet and a coface that are adjacent in a
filtration order always form an apparent pair. In the other direction, the
apparent pairs of any simplexwise filtration form a discrete gradient
\cite{bauer2021ripser}, and Bauer and Roll \cite{bauer2024wrapping} relate the
gradient pairs of a (generalised) discrete Morse function to the
zero-persistence apparent pairs of a lexicographically refined filtration
order. What we add is the consequence for the static reduction with clearing:
all column additions happen in the columns of critical cells, and on critical
rows and columns the reduction performs exactly the reduction of the Morse
complex (\cref{cor:critical-work,cor:morse-reduction}).

\paragraph{Updating persistence.}
Vineyards \cite{cohensteiner2006vines} maintain the decomposition $R = DV$ of
a filtered boundary matrix, with $V$ stored, under transpositions of adjacent
cells at linear cost per transposition. Our dynamic reduced basis (DRB) is the
same kind of object, a reduced matrix $R = BV$ with $V$ invertible and
tracked, maintained under a different set of moves. Zigzag persistence
\cite{carlsson2010zigzag} handles both insertions and deletions. Carlsson, de
Silva and Morozov \cite{carlsson2009zigzag} gave an incremental algorithm;
Milosavljevi\'c, Morozov and \v{S}kraba \cite{milosavljevic2011zigzag} bounded
the cost of $n$ additions and deletions by matrix-multiplication time; Maria
and Schreiber \cite{maria2024discrete} use discrete Morse reductions as
preprocessing, with cost depending mostly on the number of critical cells; and
FastZigzag \cite{dey2022fast} converts a simplex-wise zigzag filtration into
an ordinary filtration of a $\Delta$-complex of the same length, so that
standard persistence software applies. On graphs, zigzag persistence takes
near-linear time using dynamic connectivity structures \cite{dey2021zigzag},
and vineyard updates extend to zigzag filtrations, including expansions and
contractions \cite{dey2024vineyard}. Dey, Fan and Wang
\cite{dey2014simplicial} maintain $\mathbb{Z}_2$ simplex annotations under
elementary inclusions and vertex collapses to compute persistence along
simplicial maps. Closest to our deletion step are methods that warm-start from
an existing $R = DV$ factorisation when a filtration changes. Luo and Nelson
\cite{luo2023warmstart} extend the vineyard updates to cell insertions and
deletions processed in a batch, with cost scaling with the number of
elementary changes. Giunti and Lazovskis \cite{giunti2026pruning} update $R$,
$V$ and the representative cycles when a simplex and its cofaces are removed,
compatibly with the twist (clearing) optimisation.

\paragraph{Incremental and dynamic Betti numbers.}
Delfinado and Edelsbrunner \cite{delfinado1995incremental} compute the Betti
numbers of a complex in the 3-sphere built by adding one simplex at a time,
in $O(n\,\alpha(n))$ time. Each new simplex either creates a cycle or destroys
one (the positive/negative dichotomy that is also our insertion step); in
$S^3$ the test reduces to union--find on the complex and on its complement,
with $\beta_1$ following from the Euler characteristic. The algorithm is
insertion-only and depends on the embedding. Fully dynamic $\beta_0$ under
edge insertions and deletions is dynamic graph connectivity, solved in
polylogarithmic amortised time per update \cite{holm2001polylog}; for graphs
embedded on a surface, Eppstein \cite{eppstein2003dynamic} maintains
generators of the surface's fundamental group under edge insertions,
deletions and their duals. These bounds rely on special structure (a graph,
an ambient $S^3$, an embedding surface) and none of them covers $\beta_1$,
$\beta_2$ and $\beta_3$ of an arbitrary complex of dimension at most three
under insertions and deletions, which is our setting.

\paragraph{Positioning.}
We maintain, for each $k$, a reduced basis $R_k = B_k V_k$ with $V_k$
tracked, under column insertion and column deletion. This suffices because a
simplex's boundary never changes while the simplex exists, so rows appear and
disappear only while they are zero. Insertion is the usual incremental
reduction step. Deletion is where $V$ is needed: every column whose chain
$V_j$ contains the deleted simplex is corrected by a cycle (Case~A) or by
the column with the smallest low (Case~B). The procedure is not new in
substance. It is the rank-only specialisation of the deletion updates used
for zigzag persistence and for warm-started $R = DV$ factorisations
\cite{carlsson2009zigzag,luo2023warmstart,giunti2026pruning}, and we give a
self-contained proof because the localised re-reduction of version~1 omitted precisely
this basis tracking \cite{ouyang2025mmhm_v1}.

The difference lies in the output and in the freedom it buys. Zigzag
persistence yields the barcode of the entire sequence of complexes, a strictly
richer output: $\beta_k(K_t)$ is the number of $k$-dimensional intervals
containing~$t$, so an incremental zigzag algorithm \cite{carlsson2009zigzag}
already answers our query. In zigzag, vineyard and warm-start methods,
however, the cell order is the filtration itself and is fixed by the input.
We report only the current Betti numbers after each update, so the order
carries no meaning in our output; it is a computational device that we are
free to choose (the Morse order, under which gradient pairs are apparent pairs,
chosen in the hope of a sparser stored basis; its measured effect is mixed,
\cref{sec:e8}) and to replace at a rebuild. We claim no
locality theorem. Costs are stated in output-sensitive terms (reduction-chain
lengths and deletion footprints), we measure them, and we exhibit families on
which they grow linearly with the size of the complex. For $\beta_0$ alone,
and for insertion-only complexes in $S^3$, the specialised structures above
have better worst-case bounds. All statements are over $\mathbb{Z}_2$ and say
nothing about torsion.
